\chapter{SME E-Commerce Domain and Architectural Requirements}
\label{chap:analysis}

\section{Operational and Financial Characteristics of SME E-Commerce}
\label{sec:sme-domain}

Small and medium-sized e-commerce enterprises operate under fundamentally different economic and organizational parameters than large digital platforms \cite{ford2022,newman2021}. While hyperscale retailers optimize for multi-region availability across hundreds of engineers, SMEs optimize for capital efficiency, rapid feature experimentation, and operational simplicity. Table~\ref{tab:sme-profile} summarizes the structural constraints defining an SME e-commerce environment.

\begin{table}[h]
  \caption{Operational, financial, and technical constraints of SME e-commerce enterprises}
  \label{tab:sme-profile}
  \centering
  \begin{tabular}{p{3.6cm}p{5.2cm}p{5.4cm}}
    \toprule
    \textbf{Constraint Dimension} & \textbf{SME Operational Reality} & \textbf{Architectural Consequence} \\
    \midrule
    Engineering Team Size & 2--8 full-stack or backend developers & High cognitive and context-switching penalty if managing $>5$ repositories and CI/CD pipelines \\
    Monthly Cloud Budget & \$50--\$350/month (single VPS or small container node, 2--8 GiB RAM) & Baseline idle RAM directly determines minimum VM tier and monthly hosting expenditure \\
    Workload & 80--85\% read-heavy catalog browsing; 15--20\% stateful cart/checkout writes & Read latency ($p_{95}$) governs conversion rate during promotional campaigns \\
    Release Cadence & Weekly updates to discount rules, catalog attributes, and payment adapters & Cross-cutting edits requiring multi-service lockstep releases delay time-to-market \\
    DevOps & Automated build/deploy scripts without 24/7 dedicated SRE engineers & Distributed tracing, message-broker dead-letter queues, and sagas increase outage recovery time \\
    \bottomrule
  \end{tabular}
\end{table}

\section{Core Bounded Contexts of an E-Commerce Platform}
\label{sec:ecommerce-domains}

Applying Domain-Driven Design strategic decomposition \cite{evans2003,richardson2018}, a standard SME e-commerce platform comprises five core functional bounded contexts:
\begin{enumerate}[leftmargin=1.5em]
  \item \textbf{Catalog Management:} Product definitions, category hierarchies, brand metadata, pricing, and inventory counts. Exhibits a high read-to-write ratio ($>10:1$) and tolerates short-lived caching.
  \item \textbf{Shopping Basket:} Ephemeral or authenticated customer cart state, line-item quantities, and promotional coupon validation. Requires low-latency read/write operations during active browsing.
  \item \textbf{Ordering and Checkout:} Transactional order creation, stock reservation, shipping address persistence, and fulfillment state transitions. Demands strong ACID consistency to prevent overselling or duplicate charges \cite{kleppmann2017}.
  \item \textbf{Payment Processing:} Integration with external payment service providers (PSPs), webhook signature verification, authorization capture, and refund workflows.
  \item \textbf{Identity and Access Management:} Customer authentication, JWT/cookie session issuance, role-based access control (shopper vs.\ administrator), and profile management.
\end{enumerate}

\section{Infrastructure Cost Scaling Model for SMEs}
\label{sec:cost-model}

To connect container memory footprint directly to SME financial constraints, monthly compute cost $C_{\text{infra}}$ on standard cloud virtual private servers (VPS) or managed container instances scales as a step function of baseline resident memory $M_{\text{idle}}$ plus headroom $\Delta M_{\text{peak}}$:

\begin{equation}
\label{eq:infra-cost}
C_{\text{infra}}(A) = c_{\text{base}} \cdot \left\lceil \frac{M_{\text{idle}}(A) + \Delta M_{\text{peak}}(A) + M_{\text{OS}}}{M_{\text{tier}}} \right\rceil + C_{\text{managed}}(A)
\end{equation}

where $M_{\text{tier}}$ represents standard cloud RAM increments (e.g., $2\text{ GiB}$, $4\text{ GiB}$, $8\text{ GiB}$), $M_{\text{OS}}$ denotes host kernel and container daemon overhead ($\approx 512\text{ MiB}$), and $C_{\text{managed}}(A)$ captures auxiliary managed services (message brokers, service registries, distributed log aggregators) required by distributed topologies. A $300\text{ MiB}$--$400\text{ MiB}$ reduction in baseline runtime memory frequently allows an SME to operate reliably within a $2\text{ GiB}$ or $4\text{ GiB}$ instance tier rather than upgrading to an $8\text{ GiB}$ cluster.

\section{Structural Topology of the Evaluated Reference Systems}
\label{sec:topologies}

To eliminate programming-language and framework bias, this investigation evaluates Microsoft's two official .NET architectural reference applications—\texttt{eShopOnWeb} \cite{msft_eshoponweb} and \texttt{eShopOnContainers} \cite{msft_eshoponcontainers}—which implement the exact same e-commerce domain in C\# and ASP.NET Core. Figure~\ref{fig:arch-comparison} illustrates their contrasting structural topologies.

\begin{figure}[h]
  \centering
  \resizebox{0.98\linewidth}{!}{%
  \begin{tikzpicture}[
    svc/.style={rectangle, draw=black!80, thick, rounded corners=2pt, minimum width=2.4cm, minimum height=0.65cm, align=center, font=\scriptsize},
    db/.style={cylinder, shape border rotate=90, draw=black!80, thick, aspect=0.25, minimum width=1.25cm, minimum height=0.6cm, align=center, font=\scriptsize},
    box/.style={rectangle, draw=black!60, dashed, rounded corners=4pt, inner sep=7pt}
  ]
    % Monolith Side
    \node[svc, fill=blue!8, minimum width=4.4cm] (m_web) {ASP.NET Core Web / API Host};
    \node[svc, below=0.45cm of m_web, minimum width=4.4cm] (m_app) {ApplicationCore (Catalog, Basket, Order)};
    \node[svc, below=0.45cm of m_app, minimum width=4.4cm] (m_inf) {Infrastructure (EF Core Repositories)};
    \node[db, below=0.75cm of m_inf, minimum width=1.8cm] (m_db) {SQL Server DB};
    \node[box, fit=(m_web)(m_app)(m_inf), label=above:{\small\textbf{(a) Modular Monolith (eShopOnWeb)}}] (m_box) {};

    \draw[->, thick] (m_web) -- (m_app);
    \draw[->, thick] (m_app) -- (m_inf);
    \draw[->, thick] (m_inf) -- (m_db);

    % Microservices Side
    \node[svc, fill=orange!10, right=1.3cm of m_web, minimum width=4.7cm] (s_gw) {Envoy / YARP API Gateway};
    \node[svc, below left=0.55cm and -1.45cm of s_gw, minimum width=1.45cm] (s_cat) {Catalog API};
    \node[svc, right=0.18cm of s_cat, minimum width=1.45cm] (s_bas) {Basket API};
    \node[svc, right=0.18cm of s_bas, minimum width=1.45cm] (s_ord) {Ordering API};
    \node[svc, below=0.55cm of s_bas, minimum width=4.7cm, fill=gray!10] (s_bus) {RabbitMQ Event Bus};
    \node[db, below=0.65cm of s_bus, xshift=-1.6cm] (d_cat) {SQL};
    \node[db, below=0.65cm of s_bus] (d_bas) {Redis};
    \node[db, below=0.65cm of s_bus, xshift=1.6cm] (d_ord) {SQL};
    \node[box, fit=(s_gw)(s_cat)(s_ord)(s_bus), label=above:{\small\textbf{(b) Microservices (eShopOnContainers)}}] (s_box) {};

    \draw[->, thick] (s_gw) -- (s_cat);
    \draw[->, thick] (s_gw) -- (s_bas);
    \draw[->, thick] (s_gw) -- (s_ord);
    \draw[<->, thick] (s_cat) -- (s_bus);
    \draw[<->, thick] (s_bas) -- (s_bus);
    \draw[<->, thick] (s_ord) -- (s_bus);
    \draw[->, thick] (s_bus.south -| d_cat) -- (d_cat);
    \draw[->, thick] (s_bus.south) -- (d_bas);
    \draw[->, thick] (s_bus.south -| d_ord) -- (d_ord);
  \end{tikzpicture}%
  }
  \caption{Structural comparison of the evaluated .NET reference systems: (a) in-process layered modular monolith vs.\ (b) distributed microservices with API gateway and event bus.}
  \label{fig:arch-comparison}
\end{figure}

In \texttt{eShopOnWeb}, Clean Architecture separates \texttt{Web}, \texttt{ApplicationCore}, and \texttt{Infrastructure} into distinct C\# class libraries compiled into a single deployment unit backed by a unified SQL Server database. In \texttt{eShopOnContainers}, each bounded context runs as an autonomous Linux container with independent persistence (SQL Server for Catalog and Ordering, Redis for Basket) coordinated via an API gateway and RabbitMQ integration event bus.

\section{Chapter Summary}
\label{sec:ch2-summary}
This chapter defined the team, budget, and traffic constraints of SME e-commerce enterprises, decomposed the domain into five bounded contexts, formulated a step-function cloud memory cost model, and mapped the topologies of \texttt{eShopOnWeb} and \texttt{eShopOnContainers}.
